\begin{abstract}
We present \textbf{SETA} (Subject-aware Engagement via Threshold-tuned
Aggregation), a deployment-ready recipe for individual student
engagement recognition from classroom video. SETA combines a strong
frozen vision-language backbone (SigLIP-L) with class-balanced
bootstrap aggregation of logistic-regression probes and val-tuned
ordinal threshold decoding, achieving \textbf{$\kappa_q = 0.247~[0.201,
0.294]$} on the full 1{,}784-clip subject-disjoint DAiSEE test
set --- a \textbf{$+24\%$ relative improvement} over the standard
single-shot argmax linear-probe baseline ($\kappa_q = 0.199$),
\textbf{statistically separated} from it (lower CI 0.201 $>$ baseline
mean 0.199), and the strongest frozen-feature result we are aware of
for this benchmark.

Beyond the recipe itself, we contribute a systematic characterization
of \emph{why} this is the achievable ceiling. We introduce the
\textbf{identity-engagement entanglement (IDEP) curve} --- a
quantitative diagnostic showing that subject identity occupies a
high-rank redundant subspace of frozen features while engagement
variance is concentrated in only two-to-five linear directions that
overlap with identity. We further introduce the \textbf{within-subject
ranking diagnostic}, showing Spearman $\rho \approx 0.05$ between
predicted and true engagement within each test subject --- the model
captures subject-level priors but not per-clip variation. Cross-task
transfer to FER2013 reaches $\kappa = 0.55$ with the same frozen
encoders, confirming that the DAiSEE ceiling is engagement-specific
rather than encoder-generic. As a constructive negative control we
propose \textbf{DREAM}, a self-supervised neutral-anchor
personalization scheme, and show it fails for an explainable reason
that confirms the diagnosis: even the gentlest linear identity
removal disturbs the engagement subspace.

Our contributions are: (1) SETA, a strong reproducible recipe with
statistically meaningful improvement over the baseline;
(2) the IDEP entanglement curve as a novel diagnostic tool;
(3) the within-subject ranking finding as a behavioural fingerprint
of the ceiling; (4) DREAM as a constructive falsification of linear
identity erasure; (5) a comprehensive method-family decomposition
covering more than sixty variants across eight encoder/data-adaptation
families. All features, anchors, and scripts are released for
reproducibility.
\end{abstract}
